\chapter{القياسات الجدائية، فوبيني، تغيير المتغيّرات}
\label{ch:b3:product}

تصير نظرية لوبيغ ذات البُعد الواحد حسابًا متعدّد الأبعاد بواسطة
مبرهنتين. فتقول \emph{تونيلي--فوبيني} إن التكاملات على الجداءات
تكاملاتٌ متتالية --- أي إن التشريح مشروع، بأي ترتيب، تحت فرضيات
يمكن التحقق منها فعلًا. وتنقل \emph{صيغة تغيير المتغيّرات}
التكاملات على امتداد التماثلات التفاضلية من الصنف $\mathcal C^1$،
ويكون محدد ياكوبي هو سعر صرف الحجم؛ ونبرهن عليها كاملةً، انطلاقًا
من الحالة الخطية حيث تفسّر \emph{ما هو} المحدد. وتتوالى
التطبيقات: صيغة كعكة الطبقات، والالتفاف، \omterm{ex:b3:product:polar}{والإحداثيات القطبية}،
وحجم الكرة في البُعد $n$ --- وفي مسألة نهاية الأسبوع، صيغة
ستيرلنغ بتحليل أمين للخطأ.

\section{الجبور الجدائية من النمط $\sigma$ والقياسات الجدائية}

\begin{definition}\label{def:b3:product:sigma}
من أجل فضاءين قابلين \omterm{def:b3:measure:measure}{للقياس} $(X, \mathcal A)$ و $(Y, \mathcal B)$، يتولّد \emph{الجبر
الجدائي من النمط $\sigma$}\index{جبر جدائي سيغما@الجبر الجدائي من النمط $\sigma$}
$\mathcal A \otimes \mathcal B$ على $X \times
Y$ بواسطة \emph{المستطيلات} $A \times B$ (حيث $A \in
\mathcal A$
و $B \in \mathcal B$) --- وهي نظام من النمط $\pi$. ومن أجل $E
\subseteq X\times Y$ و $x \in X$،
تكون \emph{الشريحة} $E_x
= \{y : (x,y) \in E\}$؛ ومن أجل دالة $f$ على الجداء،
$f_x = f(x, \cdot)$.
\end{definition}

\begin{proposition}\label{prop:b3:product:sections}
(a) إذا كان $E \in \mathcal A\otimes\mathcal B$، فإن كل شريحة $E_x
\in \mathcal B$ (وبالتناظر)؛ وإذا كانت
$f$ قابلة \omterm{def:b3:measure:measure}{للقياس} بالمعنى $\mathcal
A\otimes\mathcal B$، فإن كل $f_x$ قابلة \omterm{def:b3:measure:measure}{للقياس}
بالمعنى $\mathcal
B$.
(b) $\mathcal B(\R^m)\otimes\mathcal B(\R^n) = \mathcal
B(\R^{m+n})$.
\end{proposition}

\begin{proof}
(a) المجموعات الجيدة: إن $\{E : E_x \in \mathcal B\ \forall x\}$ جبر من النمط $\sigma$ (لأن
الشرائح تتبادل مع المتممات والاتحادات القابلة للعدّ) يحتوي على
المستطيلات. وأما من أجل $f$: فإن $(f_x)^{-1}(B) = (f^{-1}(B))_x$.
(b) ($\subseteq$) مستطيلات المجموعات البوريلية: يكفي أن تكون
الصناديق مفتوحة$\times$مفتوحة بوريليةً في $\R^{m+n}$ (وهي
مفتوحة) وأن تكون المستطيلات البوريلية العامة نهاياتٍ --- وبمبدأ
المجموعات الجيدة مرة أخرى: فإن $\{A : A\times\R^n \in \mathcal B(\R^{m+n})\}$ جبر من النمط $\sigma$
يحتوي على المفتوحات؛ ثم نقاطع اثنين كهذين. ($\supseteq$) وكل
\omterm{def:b3:topology:topology}{مجموعة مفتوحة} من $\R^{m+n}$ اتحادٌ قابل للعدّ لصناديق مفتوحة ناطقة
$U\times V$: ومنه فهي محتواة في الجبر الجدائي من النمط $\sigma$.
\end{proof}

\begin{theorem}[القياس الجدائي]\label{thm:b3:product:existence}
ليكن $(X, \mathcal A, \mu)$ و $(Y, \mathcal B, \nu)$ \emph{منتهيين من النمط $\sigma$}. من أجل
كل $E \in \mathcal
A\otimes\mathcal B$، تكون الدالة $x \mapsto \nu(E_x)$ قابلة \omterm{def:b3:measure:measure}{للقياس}، ويعرّف\index{قياس
جدائي}
\[
(\mu\otimes\nu)(E) = \int_X \nu(E_x)\,\dd\mu(x)
\]
القياسَ الوحيد على $\mathcal A\otimes\mathcal B$ الذي يحقق $(\mu\otimes\nu)(A\times B) = \mu(A)\nu(B)$. وهو منتهٍ من النمط
$\sigma$، ومتناظر: إذ يُحصَل على \omterm{def:b3:measure:measure}{القياس} نفسه بمكاملة شرائح $x$
بالنسبة إلى $\nu$.
\end{theorem}

\begin{proof}
\emph{قابلية $x \mapsto \nu(E_x)$ \omterm{def:b3:measure:measure}{للقياس}.} ليكن $\nu$ منتهيًا أولًا. الصنف
$\mathcal D$ المؤلَّف من المجموعات $E$ التي يكون التطبيق من أجلها
\omterm{def:b3:lebesgue:measurable}{قابلًا للقياس} يحتوي على المستطيلات ($\nu((A\times B)_x) =
\nu(B)\mathbf 1_A(x)$) وهو نظام من النمط
$\lambda$: فمن أجل $E
\subseteq F$ في $\mathcal D$، $\nu((F\setminus E)_x) =
\nu(F_x) - \nu(E_x)$ (بالانتهاء)؛
ومن أجل $E_n \uparrow E$، $\nu((E_n)_x) \uparrow \nu(E_x)$ (بالاتصال من الأسفل)، والنهايات
الرتيبة للدوال القابلة \omterm{def:b3:measure:measure}{للقياس} قابلةٌ \omterm{def:b3:measure:measure}{للقياس}. وتشكّل المستطيلات
نظامًا من النمط $\pi$: فتعطي مبرهنة دينكين (\cref{thm:b3:measure:dynkin}) أن
$\mathcal D = \mathcal
A\otimes\mathcal B$. وإذا كان $\nu$ منتهيًا من النمط $\sigma$، كتبنا $Y =
\bigcup Y_k$
مع $Y_k \uparrow$ و $\nu(Y_k) < \infty$: فنجد $\nu(E_x) =
\lim_k\nu_k(E_x)$ حيث $\nu_k = \nu(\cdot\cap Y_k)$ منتهٍ.

\emph{\omterm{def:b3:measure:measure}{القياس}.} تنتج الجمعية من النمط $\sigma$ للمقدار $E \mapsto
\int\nu(E_x)\dd\mu$ من
\cref{cor:b3:lebesgue:additivity} (لأن شرائح المجموعات المنفصلة منفصلة). وعلى
المستطيلات تعطي $\mu(A)\nu(B)$. \emph{الوحدانية}: يتوافق مرشّحان
على النظام من النمط $\pi$ المؤلَّف من المستطيلات؛ ويعطي الانتهاء
من النمط $\sigma$ مستطيلات $X_k\times
Y_k \uparrow X\times Y$ ذات \omterm{def:b3:measure:measure}{قياس} منتهٍ: \cref{thm:b3:measure:uniqueness}.
وأما التناظر: فالبناء بالترتيب الآخر قياسٌ أيضًا يتوافق على
المستطيلات --- وهو وحيد، ومنه فهو \omterm{def:b3:measure:measure}{القياس} نفسه.
\end{proof}

\begin{definition}\label{def:b3:product:lebesgue}
\emph{قياس لوبيغ على $\R^d$}\index{قياس لوبيغ على $\R^d$} هو
$\lambda_d = \lambda\otimes\cdots\otimes\lambda$ ($d$ عاملًا؛ ويُتحقَّق من تجميعية البناء على الصناديق
وتنتشر بالوحدانية). وهو \omterm{def:b3:measure:measure}{القياس} البوريلي الوحيد الذي يعطي كل
صندوق $\prod\intoc{a_i}{b_i}$ حجمَه $\prod(b_i - a_i)$؛ وهو صامد بالانسحاب (لأن المنسحبات
تتوافق على الصناديق)، ومنتهٍ من النمط $\sigma$، وتام بعد إتمام
كاراتيودوري --- ونكتب $\lambda_d$ \omterm{def:b3:measure:measure}{للقياس} المتمَّم ونكامل تبعًا
لذلك.
\end{definition}

\section{تونيلي وفوبيني}

\begin{theorem}[تونيلي]\label{thm:b3:product:tonelli}
ليكن $\mu, \nu$ منتهيين من النمط $\sigma$، ولتكن $f \colon X\times Y \to [0,
+\infty]$ قابلة
\omterm{def:b3:measure:measure}{للقياس}. عندئذٍ تكون $x \mapsto \int_Y f_x\,\dd\nu$ قابلة \omterm{def:b3:measure:measure}{للقياس} و\index{مبرهنة تونيلي}
\[
\int_{X\times Y}f\,\dd(\mu\otimes\nu)
= \int_X\Bigl(\int_Y f(x,y)\,\dd\nu(y)\Bigr)\dd\mu(x)
= \int_Y\Bigl(\int_X f(x,y)\,\dd\mu(x)\Bigr)\dd\nu(y).
\]
\end{theorem}

\begin{proof}
الآلة المعيارية. فمن أجل $f = \mathbf 1_E$ هذه هي \cref{thm:b3:product:existence} (وصيغتها
المتناظرة). وبالخطية تصح من أجل $f \geq 0$ بسيطة. ومن أجل
$f \geq
0$ عامة: نأخذ دوالًّا بسيطة $s_n \nearrow f$ (\cref{thm:b3:lebesgue:approximation})؛
عندئذٍ $\int_Y(s_n)_x
\dd\nu \nearrow \int_Y f_x\dd\nu$ من أجل كل $x$ (بمبرهنة التقارب الرتيب في $Y$)،
ومنه تتقارب الأطراف اليسرى بمبرهنة التقارب الرتيب في $X$، بينما
$\int s_n\,\dd(\mu\otimes\nu) \nearrow \int f$ بمبرهنة التقارب الرتيب على الجداء.
\end{proof}

\begin{theorem}[فوبيني]\label{thm:b3:product:fubini}
ليكن $\mu, \nu$ منتهيين من النمط $\sigma$، ولتكن $f \in L^1(\mu\otimes\nu)$. عندئذٍ
تكون الشريحة $f_x$ \omterm{def:b3:lebesgue:l1}{قابلة للمكاملة} بالنسبة إلى $\nu$ من أجل $x$
في كل مكان تقريبًا بالمعنى $\mu$، وتكون الدالة المعرَّفة في كل
مكان تقريبًا $x \mapsto \int f_x\dd\nu$ \omterm{def:b3:lebesgue:l1}{قابلة للمكاملة}، ويساوي التكاملان المتتاليان
كلاهما $\int f\,\dd(\mu\otimes\nu)$.\index{مبرهنة فوبيني}
\end{theorem}

\begin{proof}
يبيّن تطبيق تونيلي على $\abs f$ أن $\varphi(x) = \int\abs{f_x}
\dd\nu$ ذات تكامل منتهٍ، ومنه
فهي منتهية في كل مكان تقريبًا: أي $f_x \in
L^1(\nu)$ من أجل $x$ في كل مكان
تقريبًا. ونشطر $f = f^+ - f^-$ (في الحالة الحقيقية؛ وبالمركّبات
في العقدية): فتحسب تونيلي كل تكامل متتالٍ للمقدارين $f^\pm$
بوصفه $\int f^\pm\dd(\mu\otimes\nu) <
\infty$، ويتكامل الفرق المعرَّف في كل مكان تقريبًا ليعطي
الفرق. وبالتناظر من أجل الترتيب الآخر.
\end{proof}

\begin{method}\label{met:b3:product:fubini}
لتبديل تكاملين (أو تكامل ومجموع، أو مجموعين): إذا كان المقدار
المكامَل \emph{غير سالب}، بدّل بحرية (تونيلي). وإلا فطبّق أولًا
تونيلي على $\abs f$ بالترتيب الأسهل تقديرًا؛ فإذا كانت النتيجة
منتهية، شرّعت فوبيني التبديلَ. ولا تتجاوز أبدًا فحص $\abs f$:
فللمقدار المكامَل في \cref{exo:b3:product:4} تكاملان متتاليان \emph{بقيمتين
مختلفتين}.
\end{method}

\begin{proposition}[كعكة الطبقات]\label{prop:b3:product:layercake}
من أجل $f \geq 0$ قابلة \omterm{def:b3:measure:measure}{للقياس} على $(X, \mathcal A, \mu)$ المنتهي من النمط
$\sigma$:\index{صيغة كعكة الطبقات}
\[
\int_X f\,\dd\mu = \int_0^{+\infty}\mu(\{f > t\})\,\dd t,
\qquad
\int_X f^p\,\dd\mu = p\int_0^{+\infty}t^{p-1}\mu(\{f >
t\})\,\dd t \quad (p \geq 1).
\]
\end{proposition}

\begin{proof}
نطبّق تونيلي على $\mathbf 1_{\{(x,t) : 0 < t < f(x)\}}$ على $X
\times \intoo0{+\infty}$ (وهي قابلة \omterm{def:b3:measure:measure}{للقياس}: لأنها
$\{(x,t): f(x) - t
> 0\}\cap\{t > 0\}$، وهي تركيبة من نمط بوريل \omterm{def:b3:lebesgue:measurable}{للدالة القابلة للقياس} $(x,t)\mapsto f(x) - t$):
فتعطي المكاملة في $t$ أولًا $\int f\,\dd\mu$؛ وفي $x$ أولًا،
$\int_0^\infty\mu(f >
t)\dd t$. وأما من أجل $f^p$: فنعوّض $t = s^p$ في $\int\mu(f^p > t)
\dd t$، أي نطبّق
الصيغة الأولى على $f^p$ ونغيّر المتغيّرات في التكامل ذي البُعد
الواحد ($\{f^p > s^p\} =
\{f > s\}$).
\end{proof}

\begin{theorem}[الالتفاف على $L^1$]\label{thm:b3:product:convolution}
من أجل $f, g \in L^1(\R^d, \lambda_d)$، يتقارب التكامل
\[
(f * g)(x) = \int_{\R^d} f(x - y)\,g(y)\,\dd y
\]
تقاربًا مطلقًا من أجل $x$ في كل مكان تقريبًا، ويعرّف $f * g \in
L^1(\R^d)$ مع
$\norm{f*g}_1 \leq \norm f_1\norm g_1$، ويكون $*$ تبديليًا وتجميعيًا.\index{التفاف}
\end{theorem}

\begin{proof}
المقدار $(x, y) \mapsto f(x-y)g(y)$ قابل \omterm{def:b3:measure:measure}{للقياس} (لأن $(x,y)\mapsto x - y$ متصل؛ ثم نركّب ونضرب).
وبتونيلي:
\[
\int\!\!\int \abs{f(x-y)}\abs{g(y)}\,\dd y\,\dd x
= \int\abs{g(y)}\Bigl(\int\abs{f(x - y)}\dd x\Bigr)\dd y
= \norm f_1\norm g_1 < \infty
\]
(بصمود $\lambda_d$ بالانسحاب في التكامل الداخلي). ومنه فإن
التكامل المزدوج منتهٍ؛ وتعطي فوبيني التقاربَ المطلق في كل مكان
تقريبًا وحدَّ المعيار $\norm{f*g}_1 \leq \norm
f_1\norm g_1$. وأما التبديلية: فبالتعويض $y \mapsto x - y$
(بصمود الانسحاب والانعكاس --- وصمود الانعكاس يصح على الصناديق،
ومنه في كل مكان بالوحدانية). وأما التجميعية: فبتونيلي--فوبيني
على تكامل ثلاثي.
\end{proof}

\section{تغيير المتغيّرات}

\begin{theorem}[تغيير المتغيّرات الخطي]
\label{thm:b3:product:linearchange}
من أجل $T \in GL_d(\R)$ ومن أجل $A \in \mathcal B(\R^d)$: $\lambda_d(T(A)) = \abs{\det T}\,\lambda_d(A)$؛ ومن ثَمّ $\int f(y)\dd y = \abs{\det T}\int f(Tx)\,\dd x$
من أجل $f \geq 0$ أو من أجل \omterm{def:b3:lebesgue:l1}{دالة قابلة للمكاملة}.\index{تغيير المتغيّرات
الخطي}
\end{theorem}

\begin{proof}
\omterm{def:b3:measure:measure}{القياس} $\mu_T(A) = \lambda_d(T(A))$ قياسٌ بوريلي (لأن التماثلات الطوبولوجية تحفظ
المجموعات البوريلية، \cref{pb:b3:measure:1})، وصامد بالانسحاب ($T(A + x) = T(A) + Tx$)،
ومنتهٍ على الصندوق الواحدي: ومنه، بتمييز \omterm{def:b3:measure:lebesgueouter}{قياس لوبيغ}
(\cref{exo:b3:measure:6}، الذي يصح برهانه حرفيًا في $\R^d$ بالمكعبات
الثنائية)، يكون $\mu_T = c(T)\lambda_d$ مع $c(T)
= \lambda_d(T(\intco01^d))$. والتطبيق $T \mapsto c(T)$
ضربي ($c(ST) = c(S)c(T)$، بالتركيب)، ومنه يكفي حساب $c$ على مولِّدات
$GL_d$: أي المصفوفات الأوّلية. القطرية $\operatorname{diag}(a, 1, \dots, 1)$: ترسل المكعب
الواحدي إلى صندوق حجمه $\abs a$: ومنه $c = \abs a =
\abs\det$. وتبديل
الإحداثيات: يبدّل المكعب: ومنه $c
= 1 = \abs\det$. والانسحاب المستعرض
$T(x) = x + \alpha
x_2e_1$: تكون صورة المكعب الواحدي موشورًا مقصوصًا؛ وبتونيلي
يكون قياسه $\int\lambda_1(\text{مقطع})\dd x_2
\cdots \dd x_d = 1$، لأن كل شريحة $x_1$ فترةٌ طولها $1$: ومنه
$c = 1 = \abs{\det}$. وكل مصفوفة قابلة للقلب جداءٌ لهذه (بحذف غاوس)، وكلٌّ
من $c$ و $\abs\det$ ضربي: ومنه $c(T) = \abs{\det T}$. وتنتج صيغة التكامل
بالآلة المعيارية (بالدوال المميّزة، ثم البسيطة، ثم بمبرهنة
التقارب الرتيب).
\end{proof}

\begin{theorem}[تغيير المتغيّرات]\label{thm:b3:product:changeofvar}
لتكن $U, V \subseteq \R^d$ مفتوحتين وليكن $\Phi \colon U \to V$ تماثلًا تفاضليًا من الصنف
$\mathcal C^1$. من أجل كل \omterm{def:b3:lebesgue:measurable}{دالة قابلة للقياس} $f \colon V
\to [0, +\infty]$ (أو
$f \in L^1(V)$):\index{صيغة تغيير المتغيّرات}
\[
\int_V f(y)\,\dd y = \int_U
f\bigl(\Phi(x)\bigr)\,\abs{\det D\Phi(x)}\,\dd x .
\]
\end{theorem}

\begin{proof}
نكتب $J(x) = \abs{\det D\Phi(x)}$. وقلب البرهان هو المتراجحة
\[
\lambda_d\bigl(\Phi(A)\bigr) \leq \int_A J\,\dd\lambda_d
\qquad\text{من أجل كل مجموعة بوريلية} A \subseteq U;
\tag{$*$}
\]
وترفع الخطوة 4 أدناه المتراجحةَ $(*)$ --- مطبَّقةً على كلٍّ من
$\Phi$ و $\Phi^{-1}$ --- إلى تساوي المبرهنة. ولاحظ أن
$\Phi^{-1}$ هو نفسه تماثل تفاضلي من الصنف $\mathcal C^1$ محدد
ياكوبي له $\abs{\det D\Phi^{-1}(y)} = J(\Phi^{-1}y)^{-1}$ (بقاعدة السلسلة على $\Phi\circ\Phi^{-1} = \mathrm{id}$).

\emph{الخطوة 1: المتراجحة $(*)$ من أجل المكعبات بعامل تشويه.}
لنثبّت مكعبًا مغلقًا $Q \subseteq U$ مركزه $x_0$ وضلعه $2r$ (أي
كرة بمعيار السوپريموم). الادعاء: من أجل كل $\varepsilon > 0$، إذا كان
$\Phi$ قابلًا للاشتقاق على $Q$ مع $\norm{D\Phi(x) - D\Phi(x_0)} \leq \varepsilon$ على $Q$ (\omterm{def:b3:banach:operator}{بمعيار المؤثر}
من أجل معيار السوپريموم)، فإن
\[
\Phi(Q) \subseteq \Phi(x_0) + D\Phi(x_0)\Bigl(\,\bigl(1 +
\varepsilon\norm{D\Phi(x_0)^{-1}}\bigr)\,(Q - x_0)\Bigr),
\]
لأنه من أجل $x \in Q$، تعطي متراجحة القيم المتوسطة المطبَّقة على
$\Phi(x) - \Phi(x_0) - D\Phi(x_0)(x - x_0)$ أن $\norm{\Phi(x) - \Phi(x_0) - D\Phi(x_0)(x-x_0)}_\infty \leq
\varepsilon\norm{x - x_0}_\infty \leq \varepsilon r$، ويسحب $D\Phi(x_0)^{-1}$ هذا القصورَ إلى توسيع للمكعب
بمقدار $\varepsilon\norm{D\Phi(x_0)^{-1}}\,r$. وحسب \cref{thm:b3:product:linearchange}،
\[
\lambda_d(\Phi(Q)) \leq \abs{\det D\Phi(x_0)}\,
\bigl(1 + \varepsilon\,C\bigr)^{d}\,\lambda_d(Q),
\qquad C = \sup_{Q}\norm{D\Phi(\cdot)^{-1}} .
\]

\emph{الخطوة 2: المتراجحة $(*)$ من أجل المكعبات المتراصة،
بالتجزئة.} ليكن $Q
\subseteq U$ مكعبًا متراصًّا وليكن $\varepsilon > 0$. وعلى $Q$
يكون $D\Phi$ \omterm{def:b3:topology:continuity}{متصلًا} بانتظام ويكون $\norm{D\Phi^{-1}}$ محدودًا (بالتراص)؛
نجزّئ $Q$ إلى $2^{kd}$ مكعبًا جزئيًا $Q_i$ صغيرًا بما يكفي ليكون
تذبذب $D\Phi$ على كلٍّ منها $\leq \varepsilon$. وبالخطوة 1 على كل مكعب
جزئي (مركزه $x_i$):
\[
\lambda_d(\Phi(Q)) \leq \sum_i\lambda_d(\Phi(Q_i))
\leq (1 + C\varepsilon)^d \sum_i \abs{\det
D\Phi(x_i)}\,\lambda_d(Q_i)
\leq (1 + C\varepsilon)^d\Bigl(\int_Q J + \varepsilon'\Bigr),
\]
والخطوة الأخيرة لأن $\sum_i\abs{\det D\Phi(x_i)}\mathbf
1_{Q_i} \to J$ بانتظام على $Q$ (باتصال
$\det D\Phi$) --- أي بمقارنة مجاميع ريمان. ونجعل $\varepsilon \to 0$: فتصح
$(*)$ من أجل المكعبات المتراصة.

\emph{الخطوة 3: المتراجحة $(*)$ من أجل كل مجموعة بوريلية $A$.}
دالة المجموعات $A
\mapsto \lambda_d(\Phi(A))$، على الأجزاء البوريلية من $U$، قياسٌ
(لأن $\Phi$ تقابل على $V$ يحفظ المجموعات البوريلية والانفصال
القابل للعدّ)، وكذلك $A \mapsto \int_AJ$. وكل جزء مفتوح من $U$ اتحادٌ قابل للعدّ
لمكعبات ثنائية متراصة شبه منفصلة (بالتفكيك الثنائي المعياري:
بأخذ المكعبات الثنائية الأعظمية المحتواة في المفتوحة)، والقياسان
جمعيان عليها (لأن حدود المكعبات معدومة \omterm{def:b3:measure:measure}{القياس} بالمعنى
$\lambda_d$، وصورها بالتطبيق $\Phi$ معدومة \omterm{def:b3:measure:measure}{القياس} حسب الخطوة 2
مطبَّقةً على تغطيات مكعبية رقيقة للأوجه): ومنه تنتقل $(*)$ من
المكعبات إلى المفتوحات. وأما من أجل مجموعة بوريلية عامة $A$:
فنستنفد $U$ بمتراصات $K_m \uparrow U$ مع $K_m \subseteq
\mathring K_{m+1}$، ونثبّت $m$؛
وعلى $\mathring K_{m+1}$ يكون $J$ محدودًا بعدد $M_m$ ما. وبالانتظام الخارجي
\omterm{def:b3:measure:measure}{للقياس} $\lambda_d$ (بالبرهان نفسه كما في \cref{thm:b3:measure:regularity}،
بالصناديق)، نختار مفتوحات $O_n$ تحقق $A \cap K_m \subseteq O_n
\subseteq \mathring K_{m+1}$ و $\lambda_d\bigl(O_n \setminus
(A\cap K_m)\bigr) \to 0$. عندئذٍ
\[
\lambda_d\bigl(\Phi(A\cap K_m)\bigr) \leq
\lambda_d\bigl(\Phi(O_n)\bigr) \leq \int_{O_n}J
\leq \int_{A\cap K_m}J + M_m\,\lambda_d\bigl(O_n\setminus(A\cap
K_m)\bigr) \xrightarrow[n\to\infty]{} \int_{A\cap K_m}J .
\]
ونجعل $m \to \infty$: بالاتصال من الأسفل على اليسار، وبمبرهنة
التقارب الرتيب على اليمين. وهذا يُثبت $(*)$.

\emph{الخطوة 4: التساوي وصيغة التكامل.} نمدّد أولًا $(*)$ من
المجموعات إلى التكاملات: فمن أجل كل \omterm{def:b3:lebesgue:measurable}{دالة قابلة للقياس} $g
\geq 0$
على $V$،
\[
\int_V g(y)\,\dd y \leq \int_U g(\Phi(x))\,J(x)\,\dd x .
\tag{$**$}
\]
وبالفعل، من أجل $g = \mathbf 1_B$ هذه هي $(*)$ مع $A =
\Phi^{-1}(B)$؛ وتمدّدها
الخطية إلى $g$ البسيطة، وتمدّدها مبرهنة التقارب الرتيب إلى كل
$g \geq 0$ (بالآلة المعيارية). والآن نطبّق $(**)$ مرتين: أولًا
على $g$، ثم --- من أجل التماثل التفاضلي $\Phi^{-1}$ --- على
الدالة $x \mapsto g(\Phi(x))J(x)$:
\[
\int_Vg \leq \int_U g(\Phi(x))J(x)\dd x
\leq \int_V g(y)\,J(\Phi^{-1}y)\,\abs{\det
D\Phi^{-1}(y)}\,\dd y = \int_V g ,
\]
لأن $J(\Phi^{-1}y)\abs{\det D\Phi^{-1}(y)} = \abs{\det\bigl(
D\Phi(\Phi^{-1}y)\,D\Phi^{-1}(y)\bigr)} = 1$ (بقاعدة السلسلة على $\Phi\circ\Phi^{-1} = \mathrm{id}$). فتكون جميع المتراجحات
تساويات: ومنه تصح الصيغة من أجل $g \geq 0$، ومن أجل دوال $L^1$
بالتفكيك.
\end{proof}

\begin{example}[الإحداثيات القطبية؛ الدالة الغاوسية من جديد]
\label{ex:b3:product:polar}
التطبيق $\Phi(r, \theta) = (r\cos\theta, r\sin\theta)$ تماثل تفاضلي من الصنف $\mathcal
C^1$ من $\intoo0{+\infty}\times\intoo0{2\pi}$ على
$\R^2$ منزوعًا منه نصف مستقيم (وهو معدوم \omterm{def:b3:measure:measure}{القياس})، مع
$\det D\Phi =
r$:\index{إحداثيات قطبية}
\[
\int_{\R^2}f(x, y)\,\dd x\,\dd y
= \int_0^{2\pi}\!\!\int_0^{+\infty}
f(r\cos\theta, r\sin\theta)\,r\,\dd r\,\dd\theta .
\]
ومن أجل $f = \eu^{-x^2-y^2}$، تعطي تونيلي وهذه الصيغة أن
\[
G^2 = \Bigl(\int_\R \eu^{-x^2}\dd x\Bigr)^2
= \int_{\R^2}\eu^{-x^2-y^2}
= 2\pi\int_0^\infty r\eu^{-r^2}\dd r = \pi:
\]
وهو البرهان الكلاسيكي ذو السطرين على $G = \sqrt\pi$، وقد صار
الآن مبرَّرًا تمامًا (قارنه ببرهان الوسيط في \cref{pb:b3:lebesgue:1}).
\end{example}

\begin{theorem}[حجم الكرة الواحدية]
\label{thm:b3:product:ballvolume}
لتكن $v_d = \lambda_d(B(0,1))$ في $\R^d$. عندئذٍ\index{حجم الكرة الواحدية}
\[
v_d = \frac{\pi^{d/2}}{\Gamma\bigl(\frac d2 + 1\bigr)} :
\qquad
v_1 = 2,\quad v_2 = \pi,\quad v_3 = \tfrac{4\pi}3,\quad
v_4 = \tfrac{\pi^2}2,\ \dots
\]
\end{theorem}

\begin{proof}
نحسب $I = \int_{\R^d}\eu^{-\norm x_2^2}\dd\lambda_d$ مرتين. فبتونيلي يتفكّك إلى عوامل: $I = G^d = \pi^{d/2}$. وبصيغة
كعكة الطبقات (\cref{prop:b3:product:layercake}) مع $f = \eu^{-
\norm x^2}$، ومجموعات مستوياتها
كرات: $\{f > t\} =
B\bigl(0, \sqrt{-\ln t}\bigr)$ من أجل $0 < t < 1$، وقياسها $v_d(-\ln t)^{d/2}$ (لأن التمدّد
بالعامل $\rho$ يحجّم $\lambda_d$ بالعامل $\rho^d$: \cref{thm:b3:product:linearchange})،
ومنه
\[
I = \int_0^1 v_d\,(-\ln t)^{d/2}\,\dd t
\overset{t = \eu^{-s}}{=}
v_d\int_0^\infty s^{d/2}\eu^{-s}\,\dd s
= v_d\,\Gamma\Bigl(\frac d2 + 1\Bigr).
\]
ثم نساوي. (والقيم: $\Gamma(\frac32) = \frac{\sqrt\pi}2$، $\Gamma(2) = 1$، وهكذا.) ولاحظ أن
$v_d \to 0$ حين $d \to \infty$ --- وتكمّم مسألة نهاية الأسبوع
سرعة ذلك، عبر ستيرلنغ.
\end{proof}

\section{تمارين}

\begin{exercise}[$\star$]\label{exo:b3:product:1}
ليكن $\mu$ \omterm{def:b3:measure:measure}{قياس} العدّ على $(\intcc01, \mathcal
B(\intcc01))$ (وهو ليس منتهيًا من النمط
$\sigma$) وليكن $\lambda$ \omterm{def:b3:measure:lebesgueouter}{قياس لوبيغ}، ولتكن $\Delta = \{(x,x)\}$ القطرَ في
$\intcc01^2$. برهن على أن $\Delta$ قابلة \omterm{def:b3:measure:measure}{للقياس}، واحسب
التكاملين المتتاليين للمقدار $\mathbf 1_\Delta$ بالنسبة إلى $\lambda$
و $\mu$: فهما مختلفان. وأي فرضية من فرضيات \cref{thm:b3:product:tonelli} تخفق؟
\end{exercise}

\begin{exercise}[$\star$]\label{exo:b3:product:2}
برّر التبديل واستنبط تكامل ديريكليه من جديد: من أجل $A > 0$،
\[
\int_0^A\frac{\sin x}x\,\dd x
= \int_0^A\!\!\int_0^{+\infty}\eu^{-xy}\sin x\,\dd y\,\dd x
= \int_0^{+\infty}\!\!\int_0^A \eu^{-xy}\sin x\,\dd x\,\dd y,
\]
واحسب التكامل الداخلي بصيغة مغلقة، ثم اجعل $A \to
+\infty$ (بالهيمنة
على التكامل في $y$) لتحصل على $\int_0^\infty\frac{\sin x}x\dd x = \frac\pi2$.
\end{exercise}

\begin{exercise}[$\star\star$]\label{exo:b3:product:3}
(a) برهن على أنه من أجل $f \geq 0$ قابلة \omterm{def:b3:measure:measure}{للقياس} ومن أجل $\mu$
منتهٍ: $\sum_{n\geq1}\mu(\{f \geq n\}) \leq \int f\,\dd\mu \leq
\mu(X) + \sum_{n\geq1}\mu(\{f\geq n\})$: أي إن القابلية للمكاملة هي قابلية قياسات الذيل
للجمع.
(b) استنتج أن $f \in L^1(\mu)$ (حيث $\mu$ منتهٍ) إذا وفقط إذا
كان $\sum_n\mu(\abs f \geq n) < \infty$.
\end{exercise}

\begin{exercise}[$\star\star$]\label{exo:b3:product:4}
من أجل $f(x, y) = \dfrac{x^2 - y^2}{(x^2 + y^2)^2}$ على $\intoo01^2$، برهن على
\[
\int_0^1\!\!\int_0^1 f\,\dd y\,\dd x = \frac\pi4,
\qquad
\int_0^1\!\!\int_0^1 f\,\dd x\,\dd y = -\frac\pi4
\]
\emph{(ولاحظ $f = \partial_y\bigl(\frac{y}{x^2+y^2}\bigr)$)}، وتحقق مباشرةً من $\int\!\int\abs f = +\infty$: فإن فرضية
القابلية للمكاملة في مبرهنة فوبيني ليست زخرفية.
\end{exercise}

\begin{exercise}[$\star\star$]\label{exo:b3:product:5}
(a) احسب $\mathbf 1_{\intcc01} * \mathbf 1_{\intcc01}$ صراحةً (وهي دالة خيمة)، والشكل العام للمقدار
$(\mathbf 1 * \mathbf 1 *
\mathbf 1)$.
(b) برهن على $\operatorname{supp}(f * g) \subseteq
\overline{\operatorname{supp}f + \operatorname{supp}g}$.
(c) برهن على أنه إذا كان $f \in L^1$ وكانت $g$ محدودة ومتصلة،
فإن $f * g$ متصلة. \emph{(بالتقارب المهيمن عبر اتصال الانسحاب
على $g$ المحدودة.)}
\end{exercise}

\begin{exercise}[$\star\star$]\label{exo:b3:product:6}
(a) برهن على أن المُبسَّط $\Delta_d = \{x \in \intco0\infty^d
: x_1 + \dots + x_d \leq 1\}$ حجمه $\frac1{d!}$ \emph{(بالتراجع
وفوبيني)}.
(b) استعد $v_2 = \pi$ و $v_3 = \frac{4\pi}3$ من \cref{thm:b3:product:ballvolume}، وبرهن على
$\lambda_d(\text{مجسم إهليلجي بأنصاف محاور} a_i) =
v_d\prod a_i$.
\end{exercise}

\begin{exercise}[$\star\star$]\label{exo:b3:product:7}
من أجل أي $s > 0$ يكون ما يلي منتهيًا؟ برّر \omterm{ex:b3:product:polar}{بالإحداثيات
القطبية}:
\[
\int_{B(0,1)\subseteq\R^2}\frac{\dd x\,\dd y}{(x^2 +
y^2)^{s}},
\qquad
\int_{\R^2\setminus B(0,1)}\frac{\dd x\,\dd y}{(x^2 +
y^2)^{s}} .
\]
وعمّم على $\R^d$ (بالعتبتين $s < d/2$ و $s > d/2$).
\end{exercise}

\begin{exercise}[$\star\star\star$]\label{exo:b3:product:8}
(بيتا--غاما) من أجل $p, q > 0$، لتكن $B(p, q) = \int_0^1t^{p-1}(1
- t)^{q-1}\dd t$. انطلاقًا من
$\Gamma(p)\Gamma(q)$ بوصفه تكاملًا مزدوجًا، عوّض $(x, y) = (uv,\, u(1 - v))$ (وهو تماثل تفاضلي
للربع المفتوح على $\intoo0\infty\times\intoo01$؛ واحسب محدد ياكوبي له $= u$) واخلص
إلى\index{دالة بيتا}
\[
B(p, q) = \frac{\Gamma(p)\,\Gamma(q)}{\Gamma(p + q)} .
\]
واستنتج $\int_0^{\pi/2}\sin^{2p-1}\theta\cos^{2q-1}\theta\,
\dd\theta = \frac12B(p,q)$ وقيمة تكاملات واليس $W_n = \int_0^{\pi/2}\sin^n$.
\end{exercise}

\begin{exercise}[$\star\star$]\label{exo:b3:product:9}
(صيغة النقل) لتكن $T \colon (X, \mathcal A, \mu) \to (Y,
\mathcal B)$ قابلة \omterm{def:b3:measure:measure}{للقياس} وليكن $T_*\mu(B) = \mu(T^{-1}(B))$
\emph{القياس المدفوع}\index{قياس مدفوع}. برهن على أنه من أجل كل
\omterm{def:b3:lebesgue:measurable}{دالة قابلة للقياس} $g \geq 0$ على $Y$:
\[
\int_Y g\,\dd(T_*\mu) = \int_X g\circ T\,\dd\mu
\]
(بالآلة المعيارية). ثم قارن ذلك مع \cref{thm:b3:product:linearchange}: فما المعلومة
الإضافية التي تحملها صيغة تغيير المتغيّرات ولا تحملها صيغة
النقل المجرّدة؟ \emph{(فصيغة النقل لا تعيّن $T_*\mu$ أبدًا؛ أما
مبرهنة تغيير المتغيّرات فتحسب $T_*\lambda_d$ صراحةً بوصفه
قياسًا ذا كثافة.)}
\end{exercise}

\begin{exercise}[$\star\star\star$]\label{exo:b3:product:10}
(العزوم الغاوسية) باستعمال \omterm{ex:b3:product:polar}{الإحداثيات القطبية} وفوبيني، احسب من
أجل الوزن الغاوسي المعياري على $\R^d$:
\[
\int_{\R^d}\norm x_2^2\;\eu^{-\norm x_2^2}\,\dd x
\qquad\text{و}\qquad
\int_{\R^d}x_1^2\,\eu^{-\norm x^2_2}\,\dd x,
\]
وتحقق من الاتساق ($\norm x^2 = \sum x_i^2$)، واستنتج العزم الثاني \omterm{def:b3:measure:measure}{للقياس}
$\pi^{-d/2}\eu^{-\norm x^2}\dd x$.
\end{exercise}

\begin{exercise}[$\star\star$]\label{exo:b3:product:11}
(المنحني البياني وما تحته) لتكن $f \colon \R^d \to \intco0\infty$ قابلة \omterm{def:b3:measure:measure}{للقياس}.
(a) برهن على أن \emph{ما تحت المنحني البياني} $H = \{(x, y) \in
\R^d\times\R : 0 < y < f(x)\}$ قابل
\omterm{def:b3:measure:measure}{للقياس} في $\R^{d+1}$ مع
\[
\lambda_{d+1}(H) = \int_{\R^d}f\,\dd\lambda_d :
\]
أي «التكامل هو المساحة تحت المنحني البياني»، وقد صار أخيرًا
مبرهنةً. \emph{(بالشرائح؛ وتونيلي.)}
(b) برهن على أن المنحني البياني $\{(x, f(x)) : x \in \R^d\}$ مجموعةٌ معدومة \omterm{def:b3:measure:measure}{القياس} في
$\R^{d+1}$.
(c) استنتج برهانًا من سطرين على أن الكرة $S^{d-1}$ معدومة \omterm{def:b3:measure:measure}{القياس}
بمعنى لوبيغ في $\R^d$.
\end{exercise}

\begin{exercise}[$\star\star$]\label{exo:b3:product:12}
(تكامل مزدوج شهير) باستعمال المتسلسلة الهندسية وتونيلي على
$\intoo01^2$، برهن على
\[
\int_0^1\!\!\int_0^1\frac{\dd x\,\dd y}{1 - xy}
= \sum_{n\geq1}\frac1{n^2} = \zeta(2),
\qquad
\int_0^1\!\!\int_0^1\frac{\dd x\,\dd y}{1 + xy}
= \sum_{n\geq1}\frac{(-1)^{n-1}}{n^2} = \frac{\zeta(2)}2 .
\]
(والمتطابقة الثانية للمتسلسلة: بفصل الأدلّة الزوجية عن الفردية.)
ومع $\zeta(2) = \frac{\pi^2}6$ (\cref{ch:b3:spectral})، يُقيَّم تكاملان بريئا المظهر بالقيمتين
$\frac{\pi^2}6$ و $\frac{\pi^2}{12}$؛ وأين تعمل فرضية الإيجابية في مبرهنة
تونيلي بالضبط؟
\end{exercise}

\section{مسألة: صيغة ستيرلنغ}

\begin{problem}[{مسألة نهاية الأسبوع --- $n! \sim
\sqrt{2\pi n}\,(n/\eu)^n$، بالتقارب
المهيمن}]
\label{pb:b3:product:1}
تحكم صيغة ستيرلنغ كل إحصاء مقارب في هذا الكتاب --- أحجام
الكرات، والمعاملات الثنائية، والصيغة المحلية لمبرهنة النهاية
المركزية. ونبرهن عليها من تكامل $\Gamma$ (\cref{ex:b3:lebesgue:gamma})
بواسطة \emph{طريقة لابلاس}، في صيغتها الأنقى بالتقارب المهيمن، ثم
نجمع العوائد.\index{صيغة ستيرلنغ}

\medskip
\noindent\textbf{الجزء الأول --- الصيغة.} من أجل $t > 0$،
$\Gamma(t + 1) = \int_0^\infty x^{t}\eu^{-x}\dd x$.
\begin{enumerate}
  \item عوّض $x = t + \sqrt t\,u$ وبرهن على
        \[
        \frac{\Gamma(t+1)}{t^{t}\eu^{-t}\sqrt t}
        = \int_{-\sqrt t}^{+\infty}
        \exp\Bigl(t\ln\Bigl(1 + \frac u{\sqrt
        t}\Bigr) - \sqrt t\,u\Bigr)\,\dd u
        \;=\;\int_\R g_t(u)\,\dd u,
        \]
        حيث $g_t(u) = \exp\bigl(t\ln(1 + u/\sqrt t) -
        \sqrt t\,u\bigr)\mathbf 1_{u > -\sqrt t}$.
  \item برهن على النهاية نقطةً نقطة: من أجل كل $u$ مثبَّت،
        $g_t(u)
        \to \eu^{-u^2/2}$ حين $t \to +\infty$ \emph{(انشر $\ln(1 + h)$
        حتى الرتبة الثانية)}.
  \item الهيمنة. لتكن $\varphi(h) = \ln(1 + h) - h$، بحيث يكون $g_t(u) = \exp\bigl(t\,\varphi(u/\sqrt t)\bigr)$ من أجل
        $u > -\sqrt t$. برهن على الحدّين
        \[
        \varphi(h) \leq -\frac{h^2}4 \quad (-1 < h \leq 1),
        \qquad
        \varphi(h) \leq -c\,h \quad (h \geq 1),\ \ c = 1 -
        \ln 2 > 0
        \]
        \emph{(ادرس $\varphi(h) + \frac{h^2}4$ و $\varphi(h)
        + ch$: احسب المشتقات وتحقق من
        الإشارة على كل مجال)}. واستنتج، من أجل $t \geq 1$:
        \[
        g_t(u) \leq \eu^{-u^2/4}\ \ (\abs u \leq \sqrt t),
        \qquad
        g_t(u) \leq \eu^{-cu}\ \ (u \geq \sqrt t),
        \]
        بحيث يكون $g_t(u) \leq \eu^{-u^2/4} +
        \eu^{-cu}\,\mathbf 1_{u > 0}$: أي مهيمِن قابل للمكاملة لا يتعلق
        بالوسيط $t \geq 1$.
  \item اخلص بمبرهنة التقارب المهيمن وبالتكامل الغاوسي
        (\cref{ex:b3:product:polar}):
        \[
        \Gamma(t + 1) \;\sim\;
        \sqrt{2\pi t}\;\Bigl(\frac t\eu\Bigr)^{t}
        \qquad (t \to +\infty),
        \]
        وعلى وجه الخصوص $n! \sim \sqrt{2\pi
        n}\,(n/\eu)^n$.
\end{enumerate}

\medskip
\noindent\textbf{الجزء الثاني --- العوائد.}
\begin{enumerate}[resume]
  \item (واليس) من \cref{exo:b3:product:8}، وبصيغ من نمط $W_{2n} =
        \frac\pi2\binom{2n}n4^{-n}$: استنبط
        $\binom{2n}{n} \sim \frac{4^n}{\sqrt{\pi n}}$ من ستيرلنغ، وتحقق منها مقابل العلاقة التراجعية
        $W_{n} =
        \frac{n-1}nW_{n-2}$.
  \item (أحجام الكرات تنهار) برهن على
        \[
        v_d = \frac{\pi^{d/2}}{\Gamma(\frac d2 + 1)}
        \;\sim\; \frac{1}{\sqrt{\pi d}}
        \Bigl(\frac{2\pi\eu}{d}\Bigr)^{d/2},
        \]
        ومنه فإن $v_d \to 0$ أسرع من أي متتالية هندسية؛ وجد
        البُعد الذي يعظّم $v_d$ (عدديًا: $d =
        5$).
  \item (تركّز الثنائي --- استباقًا لِما في \cref{ch:b3:clt}) باستعمال
        ستيرلنغ، برهن على التقدير المحلي، من أجل $k = n/2 + s\sqrt n/2$ مع $s$
        مثبَّت و $n$ زوجي:
        \[
        2^{-n}\binom{n}{k} \;\sim\;
        \sqrt{\frac{2}{\pi n}}\;\eu^{-s^2/2},
        \]
        أي المقطع الغاوسي المتقطّع: أي دي موافر--لابلاس في
        مهدها.
  \item وأين استعمل برهان الجزء الأول بالضبط: (أ) مبرهنة
        التقارب الرتيب أو المهيمن؛ (ب) التكامل الغاوسي؛ (ج)
        خصائص صمود \omterm{def:b3:measure:lebesgueouter}{قياس لوبيغ}؟ بجملة واحدة لكلٍّ.
\end{enumerate}

\medskip
\noindent\textbf{الجزء الثالث --- حد الخطأ: ستيرلنغ بحصر.} نضع
$d_n = \ln n! - \bigl(n + \tfrac12\bigr)\ln n + n -
\ln\sqrt{2\pi}$، بحيث يقول الجزء الأول إن $d_n \to 0$.
\begin{enumerate}[resume]
  \item برهن على $d_n - d_{n+1} = \bigl(n +
        \tfrac12\bigr)\ln\bigl(1 + \tfrac1n\bigr) - 1$.
  \item مع $t = \frac1{2n+1}$، تحقق من $\frac{n+1}n =
        \frac{1+t}{1-t}$ وانشر:
        \[
        d_n - d_{n+1} = \frac{t^2}3 + \frac{t^4}5 +
        \frac{t^6}7 + \cdots,
        \]
        واستنتج الحدّين
        \[
        \frac1{3(2n+1)^2} \;<\; d_n - d_{n+1} \;<\;
        \frac1{12n} - \frac1{12(n+1)} .
        \]
  \item تلسكب (باستعمال $d_m \to 0$) وتحقق من المتطابقة الجبرية
        اللطيفة $\frac1{3(2m+1)^2} >
        \frac1{12m+1} - \frac1{12(m+1)+1}$ من أجل $m \geq 1$، لتحصل على الحصر
        الكلاسيكي
        \[
        \sqrt{2\pi n}\Bigl(\frac n\eu\Bigr)^n
        \eu^{1/(12n+1)} \;<\; n! \;<\;
        \sqrt{2\pi n}\Bigl(\frac n\eu\Bigr)^n\eu^{1/(12n)} .
        \]
  \item نتيجتان: (a) يكون الخطأ النسبي لصيغة ستيرلنغ
        $< 10^{-6}$ متى كان $n \geq 83\,334$؛ (b) قدّر $100!$
        بأربعة أرقام معنوية يدويًا من الحصر ($100! \approx 9.3326\cdot
        10^{157}$)، وتعجّب
        قليلًا من دقة صيغة مقاربة عند $n$ منتهٍ جدًّا.
\end{enumerate}

\medskip
\noindent\textbf{الجزء الرابع --- طريق واليس: ستيرلنغ دون
الدالة الغاوسية.} تاريخيًا جاء الثابت $\sqrt{2\pi}$ من واليس لا
من غاوس؛ ويعيد هذا الجزء البرهان على ستيرلنغ مستقلًّا عن
الجزأين الأول والثاني، ومن ثَمّ يعيد البرهان على التكامل
الغاوسي. لتكن $W_n = \int_0^{\pi/2}\sin^n\theta\,
\dd\theta$.
\begin{enumerate}[resume]
  \item أثبت $W_n = \frac{n-1}nW_{n-2}$ (بالمكاملة بالتجزئة)، والصيغتين المغلقتين
        \[
        W_{2n} = \frac\pi2\binom{2n}n4^{-n}, \qquad
        W_{2n+1} = \frac{4^n}{(2n+1)\binom{2n}n},
        \]
        والمتطابقة $W_nW_{n-1} = \frac\pi{2n}$.
  \item من رتابة $(W_n)$ استنتج $W_{2n}/W_{2n+1} \to 1$، ثم
        \[
        W_{2n} \sim \frac12\sqrt{\frac\pi n}
        \qquad\text{و}\qquad
        \binom{2n}n4^{-n}\sqrt n \longrightarrow
        \frac1{\sqrt\pi} :
        \]
        أي مبرهنة واليس، محصَّلًا عليها دون ستيرلنغ.
  \item برهن، بتلسكبة الجزء الثالث وحدها (دون حاجة إلى قيمة
        الثابت)، على أن $e_n = \ln n! - (n +
        \frac12)\ln n + n$ يتقارب إلى نهاية ما $\ell$؛ وهذا
        يكافئ $n! \sim K\,n^{n+1/2}\eu^{-n}$ مع $K =
        \eu^\ell > 0$ غير معيَّن بعد.
  \item أدخل هذا التقدير المقارب في $\binom{2n}n4^{-n}\sqrt n$ وعيّن، باستعمال
        السؤال 14، القيمةَ الوحيدة الممكنة: $K = \sqrt{2\pi}$. وركّب
        المنطق: فالجزآن الثالث والرابع معًا يعطيان برهانًا
        ثانيًا كاملًا على ستيرلنغ --- ومن ثَمّ، بتشغيل تعويض
        الجزء الأول عكسيًا، تقييمًا مستقلًّا للمقدار $\int_\R\eu^{-u^2/2}\dd u = \sqrt{2\pi}$.
        ركنان، يسند كلٌّ منهما الآخر.
\end{enumerate}

\medskip
\noindent\textbf{الجزء الخامس --- عوائد أخيرة.}
\begin{enumerate}[resume]
  \item (المقطع المحلي الكامل) من أجل الأعداد الصحيحة $\abs j \leq
        K\sqrt n$
        (حيث $K$ مثبَّت)، برهن على
        \[
        \frac{\binom{2n}{n+j}}{\binom{2n}{n}} =
        \prod_{i=1}^{\abs j}\frac{n - i + 1}{n + i}
        = \exp\Bigl(-\frac{j^2}n +
        O\Bigl(\frac1{\sqrt n}\Bigr)\Bigr),
        \]
        بانتظام في $j$ \emph{(خذ اللوغاريتمات واستعمل $\ln
        \frac{1-x}{1+y} = -(x + y) + O(x^2 + y^2)$)}.
        وهذه هي الصيغة ذات الطرفين للسؤال 7 والتقدير المضبوط
        المذكور في مسألة نهاية الأسبوع من \cref{ch:b3:clt}.
  \item (استباق بواسون) برهن بستيرلنغ على $\eu^{-n}\dfrac{n^n}{n!} \sim \dfrac1{\sqrt{2\pi n}}$: أي إن منوال
        قانون بواسون ذي المتوسط الكبير $n$ يحمل الكتلة $\approx (2\pi n)^{-1/2}$،
        تمامًا كما ستتنبّأ مبرهنة النهاية المركزية.
  \item (نسب غاما) من أجل $a \in \intoo01$، برهن على $\dfrac{\Gamma(n + a)}{\Gamma(n)\,n^a} \to 1$
        باستعمال حدود ميل التحدّب اللوغاريتمي في \cref{pb:b3:lebesgue:1}
        (السؤال 14 هناك)، ومدّد إلى كل عدد حقيقي $a > 0$
        بالمعادلة الدالية. (وهذا ما تعنيه عبارة «$\Gamma(t+1) \sim$
        ستيرلنغ» بين الأعداد الصحيحة.)
  \item (الكرات، مرة أخرى) من $v_d = \pi^{d/2}/\Gamma(\frac d2
        + 1)$: كوِّن جدولًا للقيم $v_1, \dots, v_7$
        بالضبط، وتحقق من وحدانية المنوال عبر $\frac{v_d}{v_{d-2}} = \frac{2\pi}d$ (متزايد ما
        دام $d < 2\pi$، ومتناقص بعده)، وبرهن على المتطابقة
        المولِّدة اللافتة
        \[
        \sum_{k\geq0}v_{2k}\,x^{2k} = \eu^{\pi x^2} :
        \]
        أي إن جميع أحجام الكرات الواحدية ذات الأبعاد الزوجية
        محزومة في أُسّي واحد.
  \item (مقاربات الإنتروبيا) من أجل $\alpha \in \intoo01$ مثبَّت مع $\alpha n \in \N$،
        استنتج من ستيرلنغ
        \[
        \binom{n}{\alpha n} \;\sim\;
        \frac{\eu^{n\,H(\alpha)}}
        {\sqrt{2\pi\,\alpha(1-\alpha)\,n}},
        \qquad
        H(\alpha) = -\alpha\ln\alpha -
        (1-\alpha)\ln(1-\alpha) :
        \]
        أي إن معدل النمو الأُسّي للمعاملات الثنائية هو
        \emph{الإنتروبيا} $H$ --- وتحقق من أن $\alpha =
        \frac12$ يستعيد
        السؤال 5، ومن أن $H(\alpha) <
        \ln2$ من أجل $\alpha \neq \frac12$ (ومنه فإن
        المعاملات الثنائية البعيدة عن المركز مهملة أُسّيًا
        بالنسبة إلى $2^n$).
  \item (مساحات السطوح) مساحة الكرة الواحدية $S^{d-1}$ هي
        $s_{d-1} = d\,v_d$ (المبرهَن عليها بوصفها \cref{exo:b3:forms:6} في فصل الصيغ
        التفاضلية؛ ونأخذها هنا تعريفًا). كوِّن جدولًا للقيم
        $s_0, \dots, s_6$، وحدّد أكبرها ($d - 1 =
        6$، $s_6 = \frac{16\pi^3}{15} \approx 33.07$)، وبرهن على أن
        $s_{d-1} \to 0$ بسرعة تفوق الهندسية أيضًا --- فالكرات
        في الأبعاد العالية، بكل مقياس إقليدي، متلاشية الصغر.
  \item (حد التصحيح الأول) استنتج من حصر السؤال 11 أن $d_n = \frac1{12n} +
        O\bigl(\frac1{n^2}\bigr)$،
        ومنه
        \[
        n! = \sqrt{2\pi n}\,\Bigl(\frac
        n\eu\Bigr)^{n}\Bigl(1 + \frac1{12n} +
        O\Bigl(\frac1{n^2}\Bigr)\Bigr).
        \]
        وتحقق عند $n = 10$: تعطي الصيغة المجرّدة $3\,598\,696$
        (بخطأ نسبي $8.3\cdot10^{-3}$)، وتعطي المصحَّحة $3\,628\,685$
        مقابل $10! =
        3\,628\,800$ (بخطأ نسبي $3.2\cdot10^{-5}$) --- فحدٌّ واحد من
        المتسلسلة يشتري رقمين ونصف رقم.
  \item (وسيط $\Gamma$) برهن على
        \[
        \frac{1}{\Gamma(t+1)}
        \int_0^{t} x^{t}\eu^{-x}\,\dd x
        \;\longrightarrow\; \frac12
        \qquad (t \to +\infty) :
        \]
        مقاربًا، أي إن نصف كتلة المقدار المكامَل للدالة $\Gamma$
        بالضبط يقع تحت منوالها $x = t$. \emph{(شغّل تعويض الجزء
        الأول على التكامل المبتور؛ ومهيمِن السؤال 3 موجود
        أصلًا.)}
  \item (الإنتروبيا، لا مقاربًا) من أجل $\alpha \in
        \intoc0{\frac12}$ برهن على الحد،
        الصحيح من أجل \emph{كل} $n \geq 1$:
        \[
        \sum_{k=0}^{\lfloor\alpha n\rfloor}\binom nk
        \;\leq\; \eu^{n\,H(\alpha)} ,
        \]
        بمقارنة المجموع بالمقدار $\sum_k\binom
        nk\lambda^{k-\alpha n}$ من أجل الميل $\lambda =
        \frac{\alpha}{1-\alpha} \leq 1$. وتحقق
        من أن هذا الاختيار للمقدار $\lambda$ أمثلي، ووفّق بينه
        وبين السؤال 21: فالمعدل الأُسّي $H(\alpha)$ للعبارة
        المقاربة يُبلَغ بمتراجحة من سطر واحد دون أي مقاربات
        البتة.
\end{enumerate}
\end{problem}
